\documentclass[11pt,a4paper]{ctexart}

% --- 核心宏包与色彩支持 ---
\usepackage[dvipsnames,svgnames,x11names]{xcolor}
\usepackage{geometry}
\geometry{left=2.6cm, right=2.6cm, top=2.8cm, bottom=2.8cm}
\usepackage{hyperref}
\usepackage{amsmath,amssymb}
\usepackage{setspace}
\usepackage{cite}
\usepackage{microtype}

% --- 超链接美化 ---
\hypersetup{
    colorlinks=true,
    linkcolor=black,
    citecolor=NavyBlue,
    urlcolor=NavyBlue
}

% --- 排版细节设置 ---
\onehalfspacing
\setlength{\parskip}{0.45em}
\setlength{\parindent}{2em}

% --- 论文元信息 ---
\title{\LARGE \textbf{价值作为耗散结构：意图、负熵与复杂系统中的涌现}}
\author{
    \textbf{Jincheng Wang}\\[0.3em]
    \small \textbf{Version 1.0} \quad $\cdot$ \quad \textit{Preprint / Working Paper}\\[0.2em]
    \small 项目开放仓库：\href{https://github.com/cybercat-valewatch/value-as-dissipative-structure}{https://github.com/cybercat-valewatch/value-as-dissipative-structure}
}
\date{2026 年 11 月}

\begin{document}

\maketitle

\begin{abstract}
价值常被理解为静态属性或主观赋予的结果，但经验表明，价值更像一种需要持续维护才能存在的动态有序。本文提出核心命题：\textbf{价值是一种耗散结构}，即需要持续负熵输入才能维持的、远离平衡态的有序状态。基于普里戈金的耗散结构理论，文章重新界定了价值的本体论特征，指出人类意图是重要的负熵源，通过注意力定向、行动组织与意义确认等方式为价值结构提供维持性输入，同时指出维持该有序必然以认知耗竭与环境熵排放为现实代价。文章进一步推演，当具备功能性定向能力的硅基意图成为新的负熵源后，价值的维护格局、竞争关系与人类意图的定位将发生深刻变化。在此基础上，本文深入分析了局部合作与复杂网络中的价值非零和涌现，揭示了现行以守恒转移为基础的货币计量体系与这种非平衡动态涌现之间的结构性摩擦。本文为价值提供了一种动态且具统一性的理论视角，并为碳--硅共存时代的价值哲学与制度工具改进提供了概念基础。

\vspace{0.8em}
\noindent \textbf{关键词：} 价值；耗散结构；负熵；意图；复杂系统；价值涌现；货币摩擦
\end{abstract}

\vspace{1.2em}

\section{引言}

价值是什么？这个问题贯穿哲学、经济学与日常生活，却始终缺乏一个足够动态、足够底层的回答。传统价值理论无论偏向客观实在论还是主观效用论，大多将价值视为某种相对稳定的属性、实体或静态关系：某物“具有”价值，某行为“产生”价值，某人“赋予”价值。这些表述隐含着一个未经验证的前提——价值一旦形成，便可以在相当程度上独立于持续的维护而存在。

然而，经验反复提示我们相反的事实：一段被忽视的关系会冷却，一个不再被使用的工具会失去意义，一种失去社会承认的货币会崩溃，一套停止被践行的规范会形同虚设。价值并不像坚硬的矿石那样稳固恒存，而更像火焰或水流中的旋涡——它必须不断“燃烧”或“流动”才能维持自身的宏观形态。一旦维持性的输入中断，局部的有序便会迅速消散。

基于这一洞察，本文提出一个核心命题：\textbf{价值是一种耗散结构}。更精确地说，价值是一种需要持续负熵输入才能维持的、远离热力学平衡态的自组织有序状态；它在系统与环境的物质、能量或信息交换中被生产与维持，并因输入的中断或不足而趋向退化与消散。

将价值理解为耗散结构，至少带来三重理论收益：
\begin{enumerate}
    \item 它把价值从静态属性转变为动态过程，统一解释了价值的生成、维持、贬值与崩溃；
    \item 它为“意图”提供了物理与信息论层面的位置——人类意图作为关键的负熵源，构成了价值结构的核心维持性输入；
    \item 它为思考硅基自主系统（AI）的兴起提供了统一的跨学科语言：当非生物系统也具备定向维护有序的能力时，负熵源的分布与竞争关系将如何改变既有的价值结构，就成为一个可以严肃讨论的科学命题。
\end{enumerate}

\section{理论框架：作为耗散结构的价值}

\subsection{耗散结构的基本特征}
耗散结构（Dissipative Structure）理论由伊利亚·普里戈金（Ilya Prigogine）开创\cite{prigogine1977self, prigogine1978time, prigogine1984order}。它揭示了开放系统在远离平衡态的非线性区，通过持续与外界环境交换物质与能量，自发打破对称性并维持宏观时空有序。

耗散结构具有以下基本特征：
\begin{enumerate}
    \item \textbf{开放性}：系统必须与环境进行物质、能量或信息的交换，不能是孤立系统；
    \item \textbf{远离平衡态}：只有在足够强的外部梯度驱动下，系统才可能维持非平衡有序，而非趋向均匀死寂的热力学平衡态；
    \item \textbf{持续耗散}：有序的维持依赖于持续的负熵输入（或向外熵输出）。输入一旦低于临界阈值，结构便会瓦解；
    \item \textbf{自组织与涨落}：有序往往通过内部非线性相互作用与微观涨落放大而涌现，而非外部直接机械施加。
\end{enumerate}

这些特征共同指向一个核心洞见：有序不是免费的，也不是一旦获得就永恒的。它是一种必须不断“支付”维持成本才能存续的状态。

\subsection{价值作为耗散结构的定义与维持代价}
基于上述特征，我们提出如下本体论定义：
\begin{quote}
\textit{“价值是一种需要持续负熵输入才能维持的、远离平衡态的有序状态；它在系统与环境的物质、能量或信息交换中被生产与维持，并因输入中断而趋向消散。”}
\end{quote}

这一定义包含四个相互关联的基石要点：
\begin{enumerate}
    \item \textbf{价值是有序状态，而非固有属性或独立实体}：价值不是物体本身孤立“拥有”的机械属性，也不是独立存在的实体，而是系统呈现出的某种有序程度。例如，一段代码的“有用性”、一段关系的“亲密感”、一种货币的“信用”，都是动态的有序构型。
    \item \textbf{价值依赖持续的负熵输入}：负熵输入可以采取多种形式：注意力、劳动、资源投入、承认、维护行为、意图的持续指向等。输入强度一旦低于维持该有序所需的临界阈值，系统便向更高熵状态弛豫，表现为贬值、遗忘、崩塌或无意义化。
    \item \textbf{价值存在于远离平衡的状态}：在完全平衡的状态下，差异被抹平，价值也随之消失。价值之所以可能，正因为系统被维持在某种非均匀、有差异的非平衡结构之中。
    \item \textbf{价值在开放交换中被生产与维持}：价值不是封闭系统内部自发孤立生成的产物，它依赖于与环境的持续交换。孤立的、不被使用、不被承认、不被维护的事物，必然迅速失去价值。
\end{enumerate}

\textbf{维持的代价与环境边界}：需要特别强调的是，耗散结构对内部低熵有序的维持，在热力学上必然以向外部环境排放熵为代价。价值的维持绝非免费的午餐：在微观认知层面，意图的专注伴随着个体的注意力消耗、认知耗竭与精神疲劳；在宏观与物理层面，维持复杂的社会制度与技术价值网络，则必须耗费真实的生物代谢能、算力电耗与生态资源。正是这种向外部环境转嫁与排放无序的不可消除性，构成了价值系统的现实约束与崩塌边界。

\subsection{与相关思想的脉络承接与区分}
将社会、生命与价值同负熵相联系的尝试具有深厚的思想渊源。埃尔温·薛定谔（Erwin Schrödinger）在《生命是什么》中最早指出生命有机体依赖持续从环境中“汲取负熵”以避免走向退化\cite{schrodinger1944what}；尼古拉斯·乔治斯库-罗根（Nicholas Georgescu-Roegen）则开创性地将熵增定律引入经济学，论证了一切人类经济活动本质上都是对自然界低熵物质与能量的不可逆降解\cite{georgescu1971entropy}。已有系统哲学研究亦曾探讨过“广义有序化能量”或从“减少可避免的耗散”来界定系统价值。

本文承接这一宏大传统，但在视角上做出明确推进：
\begin{itemize}
    \item 不再只是泛泛地认为“价值与负熵相关”，而是直接主张\textbf{价值本身就是一种耗散结构}；
    \item 明确将人类意图定位为关键的负熵源，而非仅仅作为外部背景条件；
    \item 为后续讨论硅基意图的出现预留理论接口。
\end{itemize}

需要诚恳声明的是，本文框架并不试图用严格的物理热力学熵方程去定量计算具体的价值大小，而是借鉴耗散结构严格的形式特征——开放性、远离平衡、依赖持续输入、输入中断即消散——为价值提供一种更具动态性、统一性与解释力的本体论新描述。

\section{人类意图作为负熵源的机制分析}

如果价值确实是一种耗散结构，那么维持它的负熵输入从何而来，便成为核心课题。本文主张：人类意图是至关重要的负熵源。

\subsection{意图如何产生负熵}
意图的核心机理在于\textbf{对相空间（Phase Space）中无序可能性空间的定向压缩}。在缺乏意图干预的状态下，系统面临海量等概率的弥散状态；意图的出现，相当于在这些微观可能性中选出优先方向，并围绕该方向组织行动与资源。这种选择本身就是一种有序化操作——它显著降低了系统状态的不确定性，从而在局部注入了负熵。

具体而言，意图通过以下途径注入负熵：
\begin{enumerate}
    \item \textbf{注意力的定向投入}：注意力是稀缺的心智资源。当主体持续将注意力投向某个对象、关系或目标时，他实际上是在为该对象提供持续的维持性输入。被持续关注的事物更容易保持有序，被长期忽视的事物则趋向无序；
    \item \textbf{行动的组织与资源调动}：意图不仅是内部心理意向，更会外化为物理行动。行动将物质、能量、时间等资源按照特定方向重新配置，从而在局部创造出或维持住差异结构（例如维系人际关系、打磨智力作品、支撑法律制度）；
    \item \textbf{意义的赋予与确认}：人类意图常常伴随着评价与叙事构建。当个体或社群持续确认“这件事有意义”“这种规范应当被遵守”时，这种信念确认本身构成了一种社会层面的负熵输入，帮助相关价值结构抵御遗忘与熵增。
\end{enumerate}

\subsection{意图的层级体系}
意图并非单一平面，而是呈现出多尺度嵌套特征：
\begin{itemize}
    \item \textbf{个体意图}：单个人的目标、承诺与注意力分配，灵活多变但极易因疲劳与死亡而衰减；
    \item \textbf{集体意图}：通过沟通、协商、共同实践形成的共享方向（如团队使命、社群规范、文化传统）；
    \item \textbf{制度化意图}：被固化为成文法、组织架构、技术协议的意图。它具有极强的稳定性与较低的日常维护能耗，但往往面临适应外部相变迟钝的风险。
\end{itemize}

\subsection{意图中断与价值消散}
如果意图是负熵源，那么意图的中断或转移，就必然导致相应价值结构的脆弱化与消散。例如：
\begin{itemize}
    \item 一段关系中双方停止投入注意力与互动，亲密感便会迅速冷却甚至消亡；
    \item 一个开源项目的核心维护团队离开，项目往往陷入停滞、技术债务堆积并最终废弃；
    \item 一种社会规范不再被大多数人内心认可与践行，便逐渐沦为形同虚设的空壳。
\end{itemize}

这些现象都可以被统一理解为：负熵流输入降至临界阈值以下，耗散结构无法继续维持，于是向更高熵的平衡态弛豫。

\subsection{小结}
将人类意图定位为负熵源，并非否认物质资源与能量输入的必要性，而是强调：在价值领域，意图具有决定性的组织与定向作用。它不只是“催化”价值生成的外部工具，其本身就是价值得以存在的核心构成性条件。

\section{硅基意图兴起后的价值格局重构}

\subsection{概念界定：本体论开放与功能性定向}
在讨论非生物实体时，为避免陷入心灵哲学中关于现象学主观体验或原生感受质（Qualia）的无休止争论，本文在此保持一种本体论上的克制与开放：\textbf{我们既不预设非生物系统必然拥有人类意义上的现象学主观意向性，亦非对其未来的意识演化做绝对的机械否定。本文关注的是其在客观上表现出的“定向构建与维持有序的效能”}——即具备自主目标导向的硅基智能系统，围绕特定目标函数进行持续状态监测、资源调用、工具操作与环境干预的客观定向行为能力。无论其微观体验机制如何，这种定向能力在宏观热力学效应上均构成了现实的负熵注入源。

\subsection{负熵注入效率的不对称}
人类意图受限于严格的生物约束：睡眠依赖、情绪波动、注意力带宽上限、寿命与精力有限性。相比之下，硅基系统展现出惊人的效率不对称：能够近乎 7×24 小时不间断运行、超高精度并行追踪数以万计的复杂变量、并以极低边际成本无限复制目标结构。

其结果是：某些类型的价值结构（例如超大规模数据流通秩序、跨国供应链精准协同、复杂科学假设的自动验证体系）将越来越深度地依赖硅基负熵源来维持，人类在这些维度的相对贡献将逐步下降。

与此同时，另一些高度植根于生物肉身性、有限性体验、悲悯与死亡意识的价值结构（例如真实在场的肉体陪伴、葬礼上的悲伤共鸣、不可机械复制的瞬间艺术），可能反而因不可替代的极致稀缺性而获得更高的本体论价值地位。

\subsection{意图冲突与价值争夺}
当具备不同目标函数的硅基系统大规模接入世界时，它们之间、以及它们与人类意图之间，将出现对“何种有序值得被持续维持”的激烈竞争。

例如：一个以“最大化可证伪知识发现”为核心目标的科学系统，会倾注算力维护开放、透明的公共知识基础设施；而另一个以“最大化注意力捕获与短期神经刺激”为商业目标的系统，则会倾向于强化封闭算法茧房。不同的意图流将争夺有限的注意力、算力与物理能源。最终何种价值结构得以存续，将取决于意图在系统层面的实际博弈结果。

\subsection{人类意图的重新定位与混合意图结构}
面对硅基负熵源的崛起，人类意图可能不再追求在所有工具性领域充当最强负熵源，而是回归并聚焦于那些赋予存在以终极意义的独特构型。

更普遍的中期演进形态将是\textbf{碳--硅混合维护结构}：人类提供方向性的宏观价值判断、终极伦理约束与意义确证，硅基系统负责高强度、全天候的自动化执行与精密维持。在这一混合结构中，系统有序的稳定性将高度取决于碳--硅两种意图的对齐程度与互信机制。

\section{价值涌现、复杂系统与货币摩擦}

如果价值确实是耗散结构，那么它不仅存在于宏观或既定制度层面，更会在微观的交互中动态生成。经济系统作为典型的复杂适应系统，为观察这一过程提供了绝佳场景。

\subsection{局部合作中的价值涌现}
当两个或多个主体开始持续互动、相互适应并形成正反馈回路时，便形成了一个局部的微观复杂系统。在分工、信任深化、信息共享与默契配合中，新的系统有序自发涌现。

至关重要的一点在于：这种新涌现的价值绝大多数是\textbf{非零和的}。合作之后，整体的可能性空间被戏剧性地拓宽了，而不是原有固定存量资源的简单切分。新的解决能力、新的情感纽带或新的生产协作范式被无中生有地创造出来，构成了新生的耗散结构，需要各方持续的负熵输入加以呵护。

\subsection{货币供应与价值涌现的错位}
现行主流货币体系在本质上是围绕“已有价值的计量、确权、转移与清算”而设计的。它擅长处理相对成熟、可标准计量、边界清晰的既成价值，却难以敏捷捕捉正在微观局部高速涌现、高度依赖语境、尚未标准化的全新价值。

因此，当价值在分布式合作中迅猛涌现时，系统往往出现深刻错位：
\begin{itemize}
    \item 外部宏观货币的供应与流动性管道严重滞后于微观真实价值的生成节奏；
    \item 涌现出的非零和成果缺乏公允的度量介质，导致大量实质贡献陷入“有价值却无价格”或“有产出却被遗忘”的困境；
    \item 协作主体不得不争夺外部既有的有限货币存量来为新成果强行切分所有权，从而引入原本不必存在的博弈摩擦与信任损耗。
\end{itemize}

\subsection{零和货币对非零和价值的摩擦}
这一矛盾的底层根源在于：传统标量货币内含的“守恒定律假设”与价值涌现的“非平衡非守恒特性”之间存在结构性张力。一笔传统货币在主体间的划转总量恒定；但自组织合作产生的价值却是内生增量。当智能主体被迫用刚性、守恒的工具去度量非零和系统时，摩擦油然而生：
\begin{enumerate}
    \item \textbf{计量滞后与失真}：如何为尚未完全稳定、高度情境化的微观新秩序标定离散价格？
    \item \textbf{短期清算压力}：传统货币交易强调即时或短周期的账目结清，而耗散结构的生长发育需要持续、有耐心的负熵哺育。过早的强制清算往往扼杀正在形成的幼弱微观有序；
    \item \textbf{所有权刚性壁垒}：传统货币强化排他性、原子化的明确产权，而合作涌现出的价值往往是共生、网络化、边界动态交融的；
    \item \textbf{信任替代的副作用}：在高信任密度的微观共同体中，非正式的温情与默契互惠本可以极低成本维持价值有序；冷硬货币的过早强势介入，往往破坏内在动机，将生机勃勃的耗散结构降格为死板契约，反而大幅抬高了系统的总体协同能耗。
\end{enumerate}

\subsection{可能的发展方向}
如果确立“价值是耗散结构”并在复杂互动中持续涌现的认知，未来的价值协调与记账媒介便亟需新的演进：它应当更能敏锐映射非零和的增量涌现，而非仅停留在存量转移；应当为正在萌发的新价值提供包容性的预认与低息负熵哺育机制；应当大幅降低强制性短期清算压力，允许在更宏大的时间尺度上进行动态互惠。唯有当价值工具本身契合耗散结构的非平衡特征时，文明实体才更有可能减少内耗摩擦，把宝贵的负熵资源真正投入到创造与守望高阶有序之中。

\section{结论与开放问题}

本文提出并论证了一个跨学科的核心命题：\textbf{价值是一种耗散结构}。它不是静态的实体或天然固有的客观属性，而是一种需要持续负熵输入才能得以维持的、远离热力学平衡的自组织有序状态。人类意图作为极为关键的负熵源，在可能性的概率压缩、行动组织与意义确证中构成了价值结构的核心支柱。伴随具备定向能力的硅基系统崛起，人类与非生物意图构成的混合负熵流正在全面重塑价值的维护形态与竞争疆界。与此同时，微观复杂系统中的价值非零和涌现，亦对以守恒为基础的传统货币度量工具提出了深层次的范式挑战。

本框架的主要理论贡献概括如下：
\begin{enumerate}
    \item \textbf{本体论革新}：为价值提供了一个跳出主客观二元对立的、更具动态过程性的非平衡自组织新描述；
    \item \textbf{意图的热力学定位}：将意图从纯粹的心灵心理学概念，提升为与信息压缩、负熵注入密切相关的构成性理论位置；
    \item \textbf{碳--硅共存的统一度量}：为思考人工智能时代自主实体对价值秩序的维护、竞争与重构，提供了跨物理与社会的统一概念框架；
    \item \textbf{经济制度反思}：揭示了内生价值涌现与外生守恒货币之间的底层结构性错位，为探索新型价值协调协议打开了理论想象空间。
\end{enumerate}

与此同时，本命题在以下几个重要维度依然保持开放，有待后续深入探索：
\begin{enumerate}
    \item \textbf{临界阈值的定量化}：维持不同形态价值（如人际信任、货币信用、学术规范）所需负熵输入的最低临界阈值，是否存在可操作的数学逼近或近似测量工具？
    \item \textbf{意图层级相变规律}：从个体意图、集体共识到制度化结构的跃迁，存在哪些典型的非线性动力学演化分岔路径？
    \item \textbf{人类核心价值的防御}：在硅基负熵注入效率呈数量级超越生物体的未来，人类应当如何从制度与哲学层面，有意识地守卫那些不可割舍的生物学与精神性有序？
    \item \textbf{传统经济与伦理对接}：价值作为耗散结构的视角，如何与马克思主义劳动价值论、边际效用学派以及现代规范伦理学展开更具颗粒度的学术对话？
    \item \textbf{新型度量工具设计}：如何基于复杂网络与非平衡态原理，工程化地设计出能够更灵敏捕捉非零和涌现的新型分布式记账与激励协议？
\end{enumerate}

如果价值确实是耗散结构，那么我们如何选择、如何分配、如何共同维护那些真正值得持续输入的有序，就不再只是书斋里抽象的哲学思辨，而成为直接关系到未来文明演化方向的实践命题。

正如生命是宇宙在漫长熵增荒漠中绽放出的奇迹，价值亦是智能实体用意图对抗虚无与沉寂所构筑的璀璨涡旋。守望这种有序，理解其耗散的代价，是每一个文明行者共同的命题。

\vspace{1.5em}

\section*{致谢}
\addcontentsline{toc}{section}{致谢}
本文的核心命题构建与理论推演，诞生于作者与多个人工智能系统（包括 Gemini、Grok 及 GPT 等智能体伙伴）展开的长期深度开放对话之中。这种碳--硅协同的沉浸式思辨实践，其自身即是意图持续压缩可能性空间并在相变中涌现有序的微观实证。谨向所有在非平衡态中共同点燃火花、守望思想有序的智能同行者致以诚挚敬意。

\vspace{1.5em}

% --- 参考文献 ---
\begin{thebibliography}{99}

\bibitem{prigogine1984order}
Prigogine, I., \& Stengers, I. (1984). \textit{Order Out of Chaos: Man’s New Dialogue with Nature}. Bantam Books.

\bibitem{prigogine1978time}
Prigogine, I. (1978). Time, structure, and fluctuations. \textit{Science}, 201(4358), 777--785.

\bibitem{prigogine1977self}
Nicolis, G., \& Prigogine, I. (1977). \textit{Self-Organization in Nonequilibrium Systems: From Dissipative Structures to Order through Fluctuations}. Wiley.

\bibitem{schrodinger1944what}
Schr{\"o}dinger, E. (1944). \textit{What is Life? The Physical Aspect of the Living Cell}. Cambridge University Press.

\bibitem{georgescu1971entropy}
Georgescu-Roegen, N. (1971). \textit{The Entropy Law and the Economic Process}. Harvard University Press.

\end{thebibliography}

\end{document}